% ===================================================================
\section{Conceptual Framework}\label{sec:theory}
% ===================================================================
This section sets out the data-generating process we assume for Brent and derives the
estimand and identifying assumptions from it. The factor model in
Section~\ref{sec:theory}.1 is the organising device, treating donor selection as an economic
exercise, rendering each identifying assumption transparent, and giving the contamination
bias an explicit sign. It also exposes the residual model uncertainty that motivates the
ensemble estimator, whose construction we defer to the methodology
(Section~\ref{sec:methods}).

\subsection{A factor model of Brent}
We decompose log Brent into a component spanned by global macro-financial factors common
to many assets and an oil-specific component:
\begin{equation}
p^{\text{Brent}}_t \;=\; \underbrace{\lambda' F_t}_{\text{common global factors}}
\;+\; \underbrace{o_t}_{\text{oil-specific (supply, chokepoint, geopolitics)}}
\;+\; \varepsilon_t .
\label{eq:factor}
\end{equation}
$F_t$ collects the five drivers that move Brent in normal times, namely the global
industrial cycle, the US dollar, real interest rates, global risk appetite, and China
demand, each proxied by non-oil assets. The oil-specific term $o_t$ is the object of
interest, the supply and chokepoint component a Hormuz disruption shifts. Each donor $j$
admits its own factor representation $X_{jt}=\gamma_j'F_t+u_{jt}$ with, by construction,
zero loading on $o_t$, so a clean donor is exposed to the same global factors as
Brent but is not physically routed through oil supply. A synthetic control is then a
weighted donor basket $\hat p_t=\sum_j w_j X_{jt}$, and the logic of the method follows
directly from \eqref{eq:factor}. 

If the weights reproduce Brent's factor loadings,
$\sum_j w_j\gamma_j\approx\lambda$, the basket tracks Brent through the pre-window, where
$o_t$ is small and stable. When a chokepoint shock fires, $o_t$ jumps for Brent but not
for the basket, the two decouple, and that decoupling is the estimated premium.
Donor selection is therefore the economic exercise of spanning $F_t$ while excluding
$o_t$, that is, choosing identification over fit. This is precisely why the energy
complex, namely West Texas Intermediate (WTI) crude, refined products, and
petrocurrencies, despite the highest raw correlation with Brent, is excluded by design.
It loads on the very $o_t$ we are trying to measure.

\subsection{Potential outcomes and the estimand}
Equation~\eqref{eq:factor} describes the observed price, while the causal question
concerns the price that would have prevailed without the event. Let $Y_t^N$ be Brent's
potential log-price absent the disruption and $Y_t^I$ under it, with treatment beginning
at \T{}. The treated-unit effect is $\tau_t=Y_t^I-Y_t^N$ for $t\ge$\T{}. With only one
treated unit, $Y_t^N$ is never observed after \T{}. The synthetic supplies the missing
term, $\hat Y_t^N=\sum_j w_j X_{jt}$, and the estimated effect is the gap
$\hat\tau_t=Y_t^I-\hat Y_t^N$, reported in percent as
$\text{gap}_t=100\,[\exp(Y_t^I-\hat Y_t^N)-1]$. 

We summarise the post-window by the
\emph{average treatment effect on the treated},
\begin{equation}
\widehat{\text{ATT}} \;=\; \frac{1}{|W|}\sum_{t\in[\T{},\,\T{}+W]} \text{gap}_t .
\end{equation}
We prefer the ATT to the peak or terminal gap for two reasons. It integrates over the
entire path, including any reversion, which is the object that reserve sizing, monetary
pass-through, and fiscal buffering act on, and averaging dampens single-day jumps from
liquidity or futures-roll mechanics.

\subsection{Identifying assumptions}
Because the synthetic stands in for an unobserved counterfactual, the estimate is only as
credible as the assumptions that license it, and the factor model makes each one
explicit and checkable.
\begin{enumerate}
\item \textbf{Common-factor spanning (pre-period fit).} The donor basket must reproduce
Brent's factor loadings, so that good pre-event tracking reflects shared $F_t$ rather
than coincidence. This is checked using the pre-period root-mean-square prediction error (RMSPE),
walk-forward cross-validation, and moment matching.
\item \textbf{Donor cleanliness (SUTVA / no interference).} No donor may carry the
treatment, so the loading $\gamma_j$ must not load on $o_t$, and donor $j$ must not be hit
through a channel specific to the focal event. We used an a-priori audit and
confirmed results with model-agnostic break tests at \T{}.
\item \textbf{Regime stability.} The loadings $\lambda,\gamma_j$ must hold across the
pre-window and the projection window, so that a basket which tracked Brent before \T{}
remains a valid counterfactual after it. This is checked by distance-correlation and
structural-break diagnostics, and, across events, by comparing donor weights between
regimes.
\end{enumerate}
Assumptions~1 and 3 are testable from observed data, whereas assumption~2 is the
genuinely untestable one, since donor cleanliness concerns the unobserved
$X_{jt}^N$. The next subsection shows how the factor model nonetheless lets us sign the
direction in which a violation of assumption~2 would move the estimate.
\subsection{The direction of contamination bias}
The factor representation turns the informal worry that ``the donors might be affected by
the war'' into a quantity with a definite sign. Writing $\delta_{jt}=X_{jt}-X_{jt}^N$ for
the event-induced component of donor $j$, the estimation error is, as an exact identity,
\begin{equation}
\hat\tau_t-\tau_t \;=\; Y_t^N-\hat Y_t^N \;=\; -\sum_j w_j\,\delta_{jt}.
\label{eq:bias}
\end{equation}
The bias is minus the weight-averaged donor contamination. Its sign is therefore
an empirical matter rather than an assumption. If contaminated donors are pushed
up (oil-linked terms-of-trade, inflation$\to$rates), the synthetic is pulled up
and the ATT is under-estimated, which is conservative. If they are pushed
down (demand crowd-out, risk-off), then the ATT is over-estimated.

The identity
\eqref{eq:bias} is exact, but $\delta_{jt}$ is itself unobserved, so signing the bias in
practice requires estimating each $\delta_{jt}$ against a Brent-free reference, a
deliberately naive construction whose assumptions and limits we set out, and apply, in
Section~\ref{sec:results}. What the theory delivers here is the structure, which
is which way a given contamination pushes the estimate, not a certified magnitude.

The same identity also distinguishes two regimes of contamination by timing.
\emph{Contemporaneous} contamination occurs at \T{} and is screened by the audit and
break tests. \emph{Delayed} contamination accrues over the post-window, enters through
the projection, biases the gap toward zero, and grows with horizon, so a gap that
declines as the window lengthens is its signature, and the short-horizon gap is
the least-contaminated estimate.

A final source of error is neither identification nor contamination but
specification. The counterfactual's functional form is unknown, so no single
estimator is privileged. We address this irreducible model uncertainty with an ensemble
of five estimators rather than one, a methodological choice we set out with the
estimators themselves in Section~\ref{sec:methods}.

